\documentclass[11pt,a4paper]{article}

\usepackage{amsmath}
\usepackage{mathtools}
\usepackage{fontspec}
\usepackage{unicode-math}
\setmathfont{KpMath-Regular.otf}
\setmathfont{KpMath-Bold.otf}[version=bold]
\usepackage{microtype}
\usepackage[dvipsnames]{xcolor}
\usepackage[margin=24mm, top=28mm, bottom=28mm]{geometry}

\usepackage{fancyhdr}
\pagestyle{fancy}
\fancyhf{}
\fancyhead[L]{\small\itshape The First-Order Contrast Operator}
\fancyhead[R]{\small\itshape 17 September 2026}
\fancyfoot[C]{\thepage}
\renewcommand{\headrulewidth}{0.4pt}

\setcounter{secnumdepth}{3}
\usepackage{booktabs}
\usepackage{array}
\usepackage{multirow}
\usepackage[font=small, labelfont=bf, labelsep=period]{caption}
\usepackage{setspace}
\setstretch{1.12}
\setlength{\parskip}{0.4em}
\usepackage[colorlinks=true, linkcolor=blue!60!black, urlcolor=blue!60!black]{hyperref}

\newcommand{\bI}{\mathbf{I}}
\newcommand{\bJ}{\mathbf{J}}
\newcommand{\bA}{\mathbf{A}}
\newcommand{\bC}{\mathbf{C}}
\newcommand{\bU}{\mathbf{U}}
\newcommand{\bS}{\mathbf{S}}
\newcommand{\bW}{\mathbf{W}}
\newcommand{\bT}{\mathbf{T}}
\newcommand{\bV}{\mathbf{V}}
\newcommand{\bK}{\mathbf{K}}
\newcommand{\bq}{\mathbf{q}}
\newcommand{\bd}{\mathbf{d}}
\newcommand{\bv}{\mathbf{v}}
\newcommand{\bu}{\mathbf{u}}
\newcommand{\bta}{\boldsymbol{\tau}}
\newcommand{\bG}{\mathbf{G}}
\newcommand{\bP}{\mathbf{P}}
\newcommand{\Aop}{\mathcal{A}}
\newcommand{\dA}{\Delta\mathcal{A}}
\newcommand{\dC}{\Delta\mathbf{C}}
\newcommand{\dCeff}{\Delta\mathbf{C}_{\mathrm{eff}}}
\newcommand{\Jsix}{\mathbf{J}_{6}}
\newcommand{\Oh}{O_{h}}

\title{\bfseries The First-Order Contrast Operator\\[4pt]
\large The elastodynamic matrix--vector wave equation as a route to the\\
single-site $T$-matrix: derivation, validation, and what it actually buys}
\author{}
\date{17 September 2026}

\begin{document}
\maketitle

\begin{abstract}
\noindent
The elastodynamic matrix--vector wave equation $\partial_3\bq = \Aop\bq + \bd$
expresses the wave field as a first-order recursion in depth.  This note asks
whether the cube moment method can be applied to it directly, to obtain a
single-site $T$-matrix without passing through the Navier equation.  The answer
is yes, and the construction is not new: the thesis already carries it in
$2\tfrac12$-D.  What is new here is the $3$-D Cartesian operator, and a
measurement of what the route is worth.

Three results.  \emph{First}, the operator matrix $\Aop$ is derived from the
stiffness tensor and validated against two independent objects that agree under
one constant similarity transform; the thesis shorthand $\zeta,\chi$ are shown
to be the paper's $\nu_1,\nu_2$.  \emph{Second}, the test space is
\emph{derived} rather than chosen: $\Jsix$ is antisymmetric and
$\dA^{t}\Jsix = -\Jsix\dA$, so $\Jsix\dA$ is symmetric and the only natural
projection is $\text{test} = \Jsix\times\text{trial}$.  This dissolves an
apparent mismatch of ten conditions against nine unknowns.  \emph{Third}, and
least expected, the resulting coupling agrees with the ordinary Navier weak form
at Born order but differs at $O(\Delta c^{2})$ entirely within the traction
channel --- and that difference splits into a genuine $\Oh$-invariant
resummation \emph{and} a spurious preference for the $x_3$ axis, in comparable
amounts.  A single-site score against an independent arbiter in a homogeneous
background would mix the two and settle neither.  The directional preference
ceases to be an artefact once the background is layered, which is where the
route was wanted.

The $T$-matrix itself is assembled in \S\ref{sec:tmatrix} and the comparison is
reported channel by channel.  The sharpest statement is not a norm: $e_{12}$ and
$e_{13}$ are equivalent under $\Oh$ and so must amplify identically, and in the
first-order $T$ they split by $8.4\times10^{-4}$ while the ordinary route splits
them by exactly zero.  That split is an arbiter-free \emph{lower bound} on the
first-order error, $\ge\delta/2$; it bounds nothing about the ordinary route,
which can be symmetrically wrong, so no ordering of the two follows from it.

Every claim below is backed by an executable check; the scripts report $25/25$,
$44/44$ and $12/12$, and a separate gate binds the numerical propagator to the
derived $\Aop$.  Section~\ref{sec:open} states plainly what is \emph{not} yet
established.
\end{abstract}

\tableofcontents

%=====================================================================
\section{Conventions and objects}

\subsection{Indices, time, and units}

The positive $x_3$ axis points downward; $x_1,x_2$ are lateral.  Greek indices
$\alpha,\beta$ run over the lateral pair $\{1,2\}$ only, Latin indices over
$\{1,2,3\}$.  The time convention is $e^{-i\omega t}$ throughout, so
$\partial_t \to -i\omega$ and the particle velocity is $\bv = -i\omega\bu$.  The
lateral transform is $e^{+i(k_1x_1+k_2x_2)}$, hence $\partial_\alpha \to
ik_\alpha$.

Numerical work uses seismic units: km, km\,s$^{-1}$, g\,cm$^{-3}$, GPa.  These
are consistent, since $\text{GPa} = \text{g\,cm}^{-3}\cdot(\text{km\,s}^{-1})^2$.

Two orderings appear in the sources and must not be conflated.  The paper orders
components $(1,2,3)$; the thesis orders them $(z,x,y) = (x_3,x_1,x_2)$.  The map
between them is written explicitly in \S\ref{sec:arbiterB}.

\subsection{The field vector}

The matrix--vector wave equation is
\begin{equation}\label{eq:mvwe}
  \partial_3\bq = \Aop\,\bq + \bd,
  \qquad
  \bq = \begin{pmatrix}\bq_1\\ \bq_2\end{pmatrix},
  \quad
  \Aop = \begin{pmatrix}\Aop_{11} & \Aop_{12}\\ \Aop_{21} & \Aop_{22}\end{pmatrix},
\end{equation}
with the elastodynamic choice
\begin{equation}\label{eq:qdef}
  \bq_1 = -\bta_3, \qquad \bq_2 = \bv,
  \qquad (\bta_j)_i = \tau_{ij},
\end{equation}
so that $\bq = (-\tau_{13},\,-\tau_{23},\,-\tau_{33},\,v_1,\,v_2,\,v_3)^{t}$.
This is not an arbitrary pairing: $-\bta_3$ and $\bv$ are conjugate in the
vertical power flux, $j = \tfrac14(-\bta_3^{\dagger}\bv - \bv^{\dagger}\bta_3)$,
and that is what makes the metric of \S\ref{sec:metric} the natural one.

\subsection{The metric \texorpdfstring{$\Jsix$}{J6}}\label{sec:metric}

The symplectic metric is
\begin{equation}\label{eq:J6}
  \Jsix = \begin{pmatrix}\mathbf{0} & \bI_3\\ -\bI_3 & \mathbf{0}\end{pmatrix},
  \qquad \Jsix^{t} = -\Jsix,
  \qquad \Jsix^{2} = -\bI_6 .
\end{equation}

\paragraph{A naming warning.}  This project, following the thesis, calls this
matrix $\Jsix$; it is already defined as such in the package.  The paper calls
the \emph{same} matrix $\mathbf{N}$, and uses $\mathbf{J}$ for
$\operatorname{diag}(\bI_3,-\bI_3)$, a different object.  Translate on the way
in.  Every occurrence of $\mathbf{N}$ in the paper's equations (16)--(30)
corresponds to $\Jsix$ here.

The operator matrix is \emph{quasi-Hamiltonian} with respect to it,
\begin{equation}\label{eq:quasiham}
  \Aop^{t}\Jsix = -\Jsix\Aop,
\end{equation}
which is the paper's eq.~(27) and the thesis's \texttt{ATdef}.  Everything in
\S\ref{sec:testspace} follows from \eqref{eq:quasiham} together with the
antisymmetry of $\Jsix$, and from nothing else.

%=====================================================================
\section{The first-order system for an isotropic medium}

\subsection{From the stiffness tensor to the stiffness matrices}

For an isotropic solid,
\begin{equation}\label{eq:cijkl}
  c_{ijkl} = \lambda\delta_{ij}\delta_{kl}
           + \mu(\delta_{ik}\delta_{jl} + \delta_{il}\delta_{jk}),
  \qquad \rho_{ij} = \rho\,\delta_{ij},
\end{equation}
with $\lambda = \lambda(\mathbf{x},\omega)$, $\mu = \mu(\mathbf{x},\omega)$,
$\rho = \rho(\mathbf{x},\omega)$ and $K_c = \lambda + 2\mu$.  The stiffness
matrices are $(\bC_{jl})_{ik} = c_{ijkl}$.  In particular
\begin{equation}
  \bC_{33} = \operatorname{diag}(\mu,\mu,K_c),
  \qquad
  \bC_{33}^{-1} = \operatorname{diag}(\mu^{-1},\mu^{-1},K_c^{-1}),
\end{equation}
and $\bC_{jl}^{t} = \bC_{lj}$, which follows from $c_{ijkl} = c_{klij}$.  All
nine blocks are reproduced from \eqref{eq:cijkl} and checked entry by entry.

\subsection{Eliminating the in-plane tractions}

The components $\bta_1,\bta_2$ do not appear in \eqref{eq:qdef} and must be
removed.  Doing so introduces the reduced stiffness
\begin{equation}\label{eq:U}
  \bU_{\alpha l} = \bC_{\alpha l} - \bC_{\alpha 3}\bC_{33}^{-1}\bC_{3l},
\end{equation}
with the two structural properties $\bU_{\alpha 3} = \mathbf{0}$ and
$\bU_{\alpha\beta}^{t} = \bU_{\beta\alpha}$.  Both are checked rather than
assumed, because it is they --- not the algebra that produced them --- that make
the symmetry relations of \S\ref{sec:symrel} hold.

The isotropic $\bU_{\alpha\beta}$ is where two constants are \emph{born}:
\begin{equation}\label{eq:nu}
  \nu_1 \equiv (\bU_{11})_{11} = \frac{4\mu(\lambda+\mu)}{\lambda+2\mu},
  \qquad
  \nu_2 \equiv (\bU_{12})_{12} = \frac{2\mu\lambda}{\lambda+2\mu},
  \qquad
  \nu_1 - \nu_2 = 2\mu .
\end{equation}
They are read off $\bU$ and only then compared with the published definitions;
they are not independent quantities.  Every $\bU_{\alpha\beta}$ has a null third
row and column, which is the statement that $\tau_{33}$ has been fully
eliminated.

\subsection{The operator matrices}

The general construction gives
\begin{equation}\label{eq:Ablocks}
  \Aop_{11} = -\partial_\alpha\,\bC_{\alpha3}\bC_{33}^{-1},
  \quad
  \Aop_{12} = i\omega\rho - \tfrac{1}{i\omega}\,\partial_\alpha\bU_{\alpha\beta}\partial_\beta,
  \quad
  \Aop_{21} = i\omega\,\bC_{33}^{-1},
  \quad
  \Aop_{22} = -\bC_{33}^{-1}\bC_{3\beta}\,\partial_\beta,
\end{equation}
and for the isotropic medium these evaluate to
\begin{equation}\label{eq:A11A22}
  \Aop_{11} =
  \begin{pmatrix}
    0 & 0 & -\partial_1\tfrac{\lambda}{\lambda+2\mu}\\
    0 & 0 & -\partial_2\tfrac{\lambda}{\lambda+2\mu}\\
    -\partial_1 & -\partial_2 & 0
  \end{pmatrix},
  \qquad
  \Aop_{22} =
  \begin{pmatrix}
    0 & 0 & -\partial_1\\
    0 & 0 & -\partial_2\\
    -\tfrac{\lambda}{\lambda+2\mu}\partial_1 & -\tfrac{\lambda}{\lambda+2\mu}\partial_2 & 0
  \end{pmatrix},
\end{equation}
\begin{equation}\label{eq:A12A21}
  \Aop_{12} =
  \begin{pmatrix}
    i\omega\rho - \tfrac{1}{i\omega}(\partial_1\nu_1\partial_1 + \partial_2\mu\partial_2)
      & -\tfrac{1}{i\omega}(\partial_2\mu\partial_1 + \partial_1\nu_2\partial_2) & 0\\
    -\tfrac{1}{i\omega}(\partial_2\nu_2\partial_1 + \partial_1\mu\partial_2)
      & i\omega\rho - \tfrac{1}{i\omega}(\partial_1\mu\partial_1 + \partial_2\nu_1\partial_2) & 0\\
    0 & 0 & i\omega\rho
  \end{pmatrix},
\end{equation}
with $\Aop_{21} = i\omega\operatorname{diag}(\mu^{-1},\mu^{-1},K_c^{-1})$.

\paragraph{The ordering is the content.}  Compare the $(1,3)$ entry of
$\Aop_{11}$ with the $(3,1)$ entry of $\Aop_{22}$.  In the first the derivative
stands to the \emph{left} of the modulus and differentiates the product; in the
second it stands to the \emph{right} and differentiates the field alone.  For a
laterally invariant medium the distinction collapses and is invisible.  For a
laterally varying one it is the physics.  This is why the checks in
\S\ref{sec:arbiterA} compare \emph{operators} --- by their action on generic
fields --- rather than expressions, and why they are accompanied by a negative
control that corrupts precisely this ordering.

%=====================================================================
\section{Validating \texorpdfstring{$\Aop$}{A} against two independent objects}

\subsection{Arbiter A: the published isotropic operator matrices}\label{sec:arbiterA}

The blocks \eqref{eq:A11A22}--\eqref{eq:A12A21} are transcribed separately and
compared against those built from \eqref{eq:Ablocks}.  The comparison is made by
acting on a vector of generic symbolic functions, so it tests the derivative
ordering as well as the coefficients.  All four match.

A check that cannot fail is worth nothing, so a corrupted $\Aop_{22}$ --- with
the derivative wrongly placed outside the modulus --- is also tested, and is
rejected.

\subsection{Arbiter B: the thesis system matrix}\label{sec:arbiterB}

The thesis carries a $6\times6$ Fourier-domain system matrix $\bA(k_x,k_y)$ for
a laterally invariant medium, in a \emph{different} field vector
\begin{equation}
  \mathbf{b} = (u_3,\,u_1,\,u_2,\,\tau_{33},\,\tau_{31},\,\tau_{32})^{t},
\end{equation}
i.e.\ displacement--stress rather than velocity--traction, ordered $(z,x,y)$
rather than $(1,2,3)$, and normalised so that $\rho\omega^2$ appears where the
paper has $i\omega\rho$.  It uses the shorthand
\begin{equation}
  \gamma = \frac{\lambda}{\lambda+2\mu},\quad
  a = \frac{1}{\lambda+2\mu},\quad
  b = \frac{1}{\mu},\quad
  \zeta = \frac{4\mu(\lambda+\mu)}{\lambda+2\mu},\quad
  \chi = \frac{2\mu\lambda}{\lambda+2\mu}.
\end{equation}

\begin{quote}
\textbf{Comparing \eqref{eq:nu} with these:} $\zeta \equiv \nu_1$ and
$\chi \equiv \nu_2$.  The thesis's in-plane stiffnesses and the paper's
auxiliary moduli are the same two quantities, born in the same elimination
\eqref{eq:U}.  This is verified, not asserted.
\end{quote}

Setting $\lambda,\mu,\rho$ constant and acting on $e^{i(k_1x_1+k_2x_2)}$ reduces
$\Aop$ to a matrix.  With the constant map $\mathbf{b} = \bS\bq$ built from
$\bu = \bv/(-i\omega)$, $\bta_3 = -\bq_{1:3}$ and the permutation
$(1,2,3)\to(3,1,2)$, the identity
\begin{equation}\label{eq:similarity}
  \bA_{\text{thesis}} = \bS\,\bA_{\text{paper}}\,\bS^{-1}
\end{equation}
holds entry by entry.  Since the two matrices were derived independently and
published separately, this is evidence rather than bookkeeping.

\subsection{Structural checks}\label{sec:symrel}

Three further properties are verified on the reduced $6\times6$:
\begin{enumerate}
\item the operator symmetry relations $\Aop_{11}^{t} = -\Aop_{22}$,
      $\Aop_{12}^{t} = \Aop_{12}$, $\Aop_{21}^{t} = \Aop_{21}$, where the
      transpose is the formal adjoint under integration by parts
      ($\partial_\alpha^{t} = -\partial_\alpha$);
\item the quasi-Hamiltonian relation \eqref{eq:quasiham}, which in the
      wavenumber domain reads $\bA^{t}(-\mathbf{k})\Jsix = -\Jsix\bA(\mathbf{k})$;
\item the characteristic polynomial factorises as
      $(s^2 + k_{z,P}^2)(s^2 + k_{z,S}^2)^2$, recovering
      $\alpha^2 = (\lambda+2\mu)/\rho$ and $\beta^2 = \mu/\rho$ from a
      construction in which no wave speed was ever written down.
\end{enumerate}
Item 3 is the one that binds the algebra to physics: nothing upstream knows the
wave speeds.

%=====================================================================
\section{The contrast operator}

\subsection{Definition and the Lippmann--Schwinger equation}

Inside the scatterer the medium is $(\lambda+\Delta\lambda,\ \mu+\Delta\mu,\
\rho+\Delta\rho)$.  Writing $\Aop^{\text{tot}} = \Aop + \dA$, the total field
solves the background system driven by its own contrast,
\begin{equation}\label{eq:LS}
  (\partial_3 - \Aop)\,\bq = \dA\,\bq + \bd,
\end{equation}
which is the Lippmann--Schwinger equation of the first-order system.

\subsection{The foundation check}

Before \eqref{eq:LS} means anything, $\bq$ must actually satisfy
\eqref{eq:mvwe}.  Building $\bq$ from a \emph{generic} displacement field in a
\emph{laterally varying} medium and forming $\partial_3\bq - \Aop\bq$ gives two
halves with different strengths:
\begin{itemize}
\item rows $4$--$6$ vanish \textbf{identically}, for any $\bu$ whatever.  They
      are the constitutive relation rearranged; no equation of motion is used.
      Failure here would mean the elimination \eqref{eq:U} is wrong.
\item rows $1$--$3$ equal \textbf{exactly} the Navier body force
      $f_i = -\rho\omega^2 u_i - \partial_j\tau_{ij}$.  Anything left over would
      be a defect in $\Aop_{11}$ or $\Aop_{12}$.
\end{itemize}
Both hold.  This is a considerably stronger statement than the plane-wave and
eigenvalue checks of \S\ref{sec:symrel}, because $\bu$ is generic and the
moduli vary laterally, so every ordering term in \eqref{eq:A11A22} is exercised.

\subsection{\texorpdfstring{$\dA$}{DeltaA} is not linear in the contrast}
\label{sec:nonlinear}

The Navier potential is exactly linear in $(\Delta\lambda,\Delta\mu,\Delta\rho)$.
$\dA$ is not, because the elimination inverted $\bC_{33}$:
\begin{equation}\label{eq:resum}
  (\dA_{21})_{11} = i\omega\left(\frac{1}{\mu+\Delta\mu} - \frac{1}{\mu}\right)
  = -\,\frac{i\omega\,\Delta\mu}{\mu(\mu+\Delta\mu)},
\end{equation}
whose ratio to its linearisation $-i\omega\Delta\mu/\mu^2$ is exactly
$\mu/(\mu+\Delta\mu)$.  At the project's validated moderate contrast
($\alpha = 5$, $\beta = 3$, $\rho = 2.5$, so $\lambda = 17.5$, $\mu = 22.5$ GPa;
$\Delta\lambda = 2$, $\Delta\mu = 1$ GPa):
\begin{center}
\begin{tabular}{lr}
\toprule
quantity & value\\
\midrule
$\Delta\mu/\mu$ & $0.04444$\\
exact\,/\,linear for $\Delta b = \Delta(1/\mu)$ & $0.95745$\\
exact\,/\,linear for $\Delta a = \Delta(1/K_c)$ & $0.93985$\\
\bottomrule
\end{tabular}
\end{center}
The first-order form has already summed part of the local response exactly.
\S\ref{sec:whatitbuys} measures what that is worth, and finds the answer is not
the obvious one.

%=====================================================================
\section{The augmentation}

\subsection{Why a derivative must be moved}

Inside the scatterer the contrasts are constant, so every coefficient of $\dA$
carries the indicator $\mathbf{1}_V$.  A derivative standing to the \emph{left}
of such a coefficient differentiates the indicator and produces a surface delta
on the cube faces.  Integration by parts over all space moves that derivative
onto the test function instead, where it meets something smooth:
\begin{equation}\label{eq:byparts}
  \int \phi_i\,\bigl(-\partial_\alpha( m\,\mathbf{1}_V\,\psi_j)\bigr)\,d\mathbf{x}
  = +\int (\partial_\alpha\phi_i)\, m\,\mathbf{1}_V\,\psi_j\,d\mathbf{x},
\end{equation}
and the face term never arises.  This is the step the thesis describes as
transferring the derivatives onto the source and receiver coordinates of the
propagator, and it is what converts a singular integro-differential equation
into a nonsingular, compact integral equation.

\subsection{The augmented dimension}

After the transfer every term has the form
$(\partial^{a}\phi)^{t}\,\cdot\,\text{const}\,\cdot\,(\partial^{b}\psi)$ with
$a,b\in\{0,1\}$, so the form is \emph{multiplicative} on the enlarged vector
$\mathcal{L}\psi = (\psi,\ \partial_1\psi,\ \partial_2\psi)$ --- nominally
$6+6+6 = 18$ slots.  The true dimension is smaller and is counted, not assumed.

\begin{quote}
\textbf{Result.}  Exactly $10 = 6+4$ slots are used.  The four extra are
$\partial_1 v_1,\ \partial_2 v_1,\ \partial_1 v_2,\ \partial_2 v_2$: the lateral
gradients of the two \emph{horizontal} velocity components, and nothing else.
\end{quote}

Setting $\partial_2 \to ik_y$, as a $2\tfrac12$-D formulation does, makes two of
those four multiplicative, leaving $6+2 = 8$.  That is precisely the thesis's
count for its augmented coupling matrix, reproduced here from an independent
$3$-D Cartesian derivation.  The agreement is a corroboration of the thesis
construction, not a restatement of it.

The transfer identity is verified three ways: the augmented matrix is constant
(no derivative operators survive); the augmented form reproduces the transferred
bilinear; and the untransferred and transferred bilinears differ by an
explicitly exhibited lateral divergence $\partial_\alpha P_\alpha$.

%=====================================================================
\section{The test space, derived}\label{sec:testspace}

\subsection{The argument}

At this point there are ten augmented slots and, for an affine trial space, nine
parameters.  The temptation is to choose nine conditions.  That temptation must
be resisted: a subset chosen by scoring against the arbiter it is later judged
by is a fitted parameter wearing a derivation's clothes.

There is no choice to make.  $\Jsix$ is \emph{anti}symmetric, and
\eqref{eq:quasiham} applies to $\dA$ as well as to $\Aop$ (it is linear in the
operator), so
\begin{equation}\label{eq:symmetric}
  (\Jsix\dA)^{t} = \dA^{t}\Jsix^{t} = \dA^{t}(-\Jsix) = \Jsix\dA .
\end{equation}
$\Jsix\dA$ is a \textbf{symmetric} matrix, so
\begin{equation}
  b[\phi,\psi] = \phi^{t}\,\Jsix\dA\,\psi
\end{equation}
is a symmetric bilinear form.  A symmetric form admits exactly one natural
projection:
\begin{equation}\label{eq:testspace}
  \boxed{\ \text{test space} \;=\; \Jsix \times \text{trial space}. \ }
\end{equation}
Verified, together with the negative control: drop $\Jsix$ and the symmetry is
lost, both at operator level and in the projected system.

\subsection{The ten-versus-nine mismatch dissolves}

The ten slots were never test conditions.  They measure how many derivative
channels the \emph{form} touches --- its representation --- not how many
independent functionals it supplies.  Each trial function $\psi_k$ gives exactly
one test functional $(\Jsix\psi_k)^{t}$, so the number of conditions equals the
trial dimension by construction.

\subsection{The effective coupling matrix}

For the affine trial space --- nine parameters, three rigid translations and six
uniform strains --- the projected coupling
\begin{equation}\label{eq:dCeff}
  (\dCeff)_{kl} = \int_V (\mathcal{L}\psi_k)^{t}\,\dC\,(\mathcal{L}\psi_l)\,dV
\end{equation}
is the $3$-D Cartesian analogue of the thesis's effective coupling matrix.  It
is $9\times9$, symmetric, of full rank $9$, vanishes with the contrast, and
--- because the trial functions are polynomials --- requires \textbf{no Green's
function at all}: the cube integral is exact.

The rigid-translation block is diagonal,
\begin{equation}\label{eq:rigid}
  (\dCeff)_{1:3,1:3} = i\omega^{3}\,\Delta\rho\,V\,\bI_3,
  \qquad V = (2a)^3 .
\end{equation}

\paragraph{The power of $\omega$ is a basis statement, not a physical one.}
$\dA$'s density entry is $i\omega\Delta\rho$ acting on the velocity slot; the
trial carries $\bv = -i\omega\bu$ there; and $\Jsix$ pulls the test function's
velocity slot into the traction rows, supplying a second $-i\omega$.  Hence
$i\omega\cdot(-i\omega)\cdot(-i\omega) = i\omega^{3}$.  The displacement-basis
answer $i\omega\Delta\rho V$ is wrong by $\omega^2$.  This is the
displacement-versus-flux confusion in a new place, and the check asserts the
correct factor \emph{and} asserts the incorrect one is incorrect, so it cannot
drift back.

%=====================================================================
\section{What the first-order coupling actually buys}\label{sec:whatitbuys}

\subsection{The trial spaces coincide}

It is tempting to attribute any difference from the Navier formulation to the
nonlinearity of \S\ref{sec:nonlinear}, or to a different ansatz.  Neither is
right.  Inside the scatterer the moduli are constant, so
$\bta_3 = \bC_{33}^{\text{tot}}\,\boldsymbol\varepsilon_{i3}$ is a constant
\emph{invertible} map.  A uniform strain gives a uniform $\bta_3$ and an affine
$\bv$, and conversely: ``$\boldsymbol\varepsilon$ uniform'' and ``$\bta_3$
uniform plus in-plane $\boldsymbol\varepsilon$ uniform'' describe the same
nine-parameter family.

The two formulations are therefore \textbf{different projections of the same
residual onto the same trial space}.  The difference lives in the test
conditions.

\subsection{The density channel cannot differ}

With $\Delta\lambda = \Delta\mu = 0$, $\dA$ is exactly $i\omega\Delta\rho$ in
rows $1$--$3$, columns $4$--$6$, and reaches \emph{no} part of the stress-glut
slot: a density contrast is a pure body force.  The propagator's body-force
velocity block is $-i\omega\bG$, so both routes return
\begin{equation}
  A_u = \frac{1}{1 - \omega^2\Delta\rho\,\Gamma_0},
  \qquad \Gamma_0 = \int_V \bG\,dV,
\end{equation}
identically.  The modulus channels are the only place the two can part.

\subsection{Born agreement, and an \texorpdfstring{$O(\Delta c^2)$}{O(Dc2)} difference}

Reading a proportionality constant off the rigid block alone gives $i\omega$,
and testing it on the remaining eighty entries \textbf{refutes} the hypothesis
$\dCeff = i\omega\,\dC_{\text{Navier}}$.  The two couplings agree exactly at
$O(\Delta c^{0})$ and $O(\Delta c^{1})$ --- identical Born limits, so nothing is
broken --- and differ from second order.

\begin{quote}
\textbf{Where.}  Eleven of eighty-one entries differ.  With the parameter order
$(c_1,c_2,c_3,\ e_{11},e_{22},e_{33},\ e_{12},e_{13},e_{23})$ the nonzero block
is $\{e_{11},e_{22},e_{33}\}^{2}$ together with $(e_{13},e_{13})$ and
$(e_{23},e_{23})$.
\end{quote}

That is exactly the set of strains $\bta_3$ sees:
\begin{equation}
  \tau_{13} = 2\mu\,e_{13},\qquad
  \tau_{23} = 2\mu\,e_{23},\qquad
  \tau_{33} = \lambda(e_{11}+e_{22}+e_{33}) + 2\mu\,e_{33}.
\end{equation}
The in-plane shear $e_{12}$ is \textbf{untouched}, as it must be: $\tau_{12}$ was
eliminated, not resummed.  This was a falsifiable prediction about which entries
may be nonzero, and it held.  A representative entry is
\begin{equation}
  \bigl(\dCeff - i\omega\,\dC_{\text{Navier}}\bigr)_{e_{11}e_{11}}
  = -\,\frac{i\omega V\,\Delta\lambda^{2}}{\lambda+2\mu},
\end{equation}
pure second order, divided by $K_c$: the remainder of the geometric series for
$1/(K_c+\Delta K_c)$.

\subsection{Half of it is a broken symmetry}\label{sec:oh}

An isotropic cube in an isotropic whole space has an $\Oh$-symmetric
$T$-matrix; no direction is special.  But the first-order system is a recursion
in \emph{depth}: it eliminated $\tau_1,\tau_2$ and kept $\tau_3$, so $x_3$ is
preferred by construction.  The tell is already above --- $(e_{13},e_{13})$ and
$(e_{23},e_{23})$ differ while $(e_{12},e_{12})$ does not, and those three are
$\Oh$-equivalent.

Applying the cyclic element $x\to y\to z\to x$ (positive controls: $P^3 = \bI$,
and the Navier coupling \emph{is} invariant):

\begin{center}
\begin{tabular}{lrrr}
\toprule
contrast (GPa) & invariant & breaking & ratio\\
\midrule
pure shear, $\Delta\mu = 1$                       & $2.62\times10^{-2}$ & $1.94\times10^{-2}$ & $1.3520$\\
moderate, $(2,\,1,\,0.1)$                         & $5.17\times10^{-4}$ & $2.93\times10^{-4}$ & $1.7669$\\
strong, $(20,\,10,\,0.5)$                         & $1.05\times10^{-2}$ & $5.92\times10^{-3}$ & $1.7669$\\
\bottomrule
\end{tabular}
\captionof{table}{Decomposition of $\dCeff - i\omega\dC_{\text{Navier}}$,
relative to $\|\dCeff\|$, at $\lambda = 17.5$, $\mu = 22.5$ GPa.  The two mixed
rows share a contrast \emph{direction} $(\Delta\lambda{:}\Delta\mu = 2{:}1)$ at
different magnitudes and return the identical ratio, so the split is a function
of contrast direction alone --- as a leading-order $O(\Delta c^2)$ split must be.}
\end{center}

So the difference is neither pure artefact nor pure physics.  It carries both a
genuine $\Oh$-invariant resummation the Navier form lacks, and a spurious
preference for $x_3$.

\paragraph{A caveat on this test.}  The group used is the cyclic subgroup of
order $3$, not all $48$ elements of $\Oh$.  Invariance under it is necessary for
$\Oh$ invariance, not sufficient.

\subsection{What follows, and what does not}

Two readings are available and the measurement above does not separate them.

\emph{Against the route:} an approximation that breaks a symmetry the exact
answer possesses has introduced an error that no propagator can remove.

\emph{For the route:} the correction is not scattered.  It lands exactly on the
traction channel, treating $\bC_{33}$ exactly instead of leaving it to the
Neumann series, and leaves the in-plane channel alone.  Vertical $P$ incidence
excites $e_{33}$; vertical $S$ incidence excites $e_{13}$ or $e_{23}$; both are
corrected.  The in-plane shear, which the side-scattered field carries, is
unchanged.  If the response is dominated by reflection and transmission, the
asymmetry is ``$z$ treated better'' rather than ``$z$ treated worse''.

\begin{quote}
\textbf{Consequence for any comparison.}  Do not score the raw coupling against
an independent arbiter in a homogeneous background.  Three effects are in play
at comparable magnitude --- the Galerkin-versus-collocation projection, the
$\Oh$-invariant resummation, and the artefact --- and a single number separates
none of them.  Any comparison must be channel-resolved.
\end{quote}

Whether the resummation is an \emph{improvement} is not established here.  The
local inversion $1/(K_c+\Delta K_c)$ is exact; whether that makes the closed
$T$-matrix better depends on how the ansatz error and the series truncation
interact, and those do not simply add.

%=====================================================================
\section{The propagator, and the bridge to it}

\subsection{It already exists}

Everything downstream needs the Green's matrix of \eqref{eq:mvwe}.  It does not
need to be built: the stratified $6\times6$ Green's function and its validated
wrapper already return it, in the basis
$(u_z,u_x,u_y,T_{zz},T_{xz},T_{yz})$ --- which is exactly the basis the thesis
writes $\bA$ in.  That the two are the same object had, however, been
\emph{assumed}.

\subsection{The bridge identity}

Away from the source plane the Green's matrix solves the homogeneous system in
its receiver coordinate, so for two receiver depths in one uniform layer
\begin{equation}\label{eq:bridge}
  \boldsymbol\Gamma(z_2,z') = \exp\!\bigl(\bA\,(z_2-z_1)\bigr)\,\boldsymbol\Gamma(z_1,z'),
\end{equation}
with no free parameters.  The identity is \emph{local}: it holds however
complicated the medium between the source and $z_1$ is, because the ODE knows
only the layer the receiver sits in.  The gate therefore places the source above
a fast slab and the receivers below it, so transmission, reflection, conversion
and interbed multiples are all present in $\boldsymbol\Gamma$ and
\eqref{eq:bridge} must still hold exactly.

\subsection{Reporting residuals against the arithmetic floor}

A bare relative residual is not comparable across incidence angles.  Beyond the
critical slowness the propagator carries both a growing and a decaying
evanescent branch, so $\exp(\bA\,dz)$ has dynamic range $\sim e^{|k|dz}$ ---
about $1.6\times10^{7}$ at the widest angle used.  Forming
$\exp(\bA\,dz)\,\boldsymbol\Gamma(z_1)$ then amplifies the rounding already
present by $\|\exp(\bA\,dz)\|$, and no correct implementation can beat
\begin{equation}
  \text{floor} = \epsilon\;\|\exp(\bA\,dz)\|\;
  \frac{\|\boldsymbol\Gamma(z_1)\|}{\|\boldsymbol\Gamma(z_2)\|}.
\end{equation}
Residuals below are quoted in units of that floor, so $O(1)$ means the identity
holds to the last bit available at that conditioning.

\begin{center}
\begin{tabular}{rrrrr}
\toprule
$k_x$ & $k_y$ & $|k|/\omega$ & raw & corrected\\
\midrule
$2.00$  & $1.30$ & $0.0759$ & $7.68\times10^{13}$ & $1.26\times10^{-2}$\\
$0.40$  & $0.05$ & $0.0128$ & $2.95\times10^{14}$ & $2.43\times10^{-2}$\\
$14.00$ & $9.00$ & $0.5298$ & $4.31\times10^{13}$ & $9.37\times10^{-3}$\\
\midrule
\multicolumn{3}{r}{control, $\zeta\leftrightarrow\chi$ swapped} & $7.68\times10^{13}$ & $1.36\times10^{12}$\\
\bottomrule
\end{tabular}
\captionof{table}{Residual of \eqref{eq:bridge} in units of the arithmetic
floor.  $\omega = 31.4159$ rad\,s$^{-1}$, $dz = 1$ km.  The third row is
evanescent for both branches.}
\end{center}

\begin{quote}
\textbf{The corrected object is the physical state vector of
$\partial_3\bq = \Aop\bq$}; the raw one misses by fourteen orders of magnitude.
\end{quote}

The negative control swaps $\zeta$ and $\chi$ --- two in-plane stiffnesses that
differ by exactly $2\mu$, a plausible transcription slip that no symmetry check
would catch.  It is rejected by fourteen orders.

%=====================================================================
\section{The single-site \texorpdfstring{$T$}{T}-matrix}\label{sec:tmatrix}

This is the object the construction was for.  On the nine-parameter affine
trial space,
\begin{equation}\label{eq:T}
  \boxed{\ \mathbf{T} = \dCeff\bigl(\bI_9 - \mathcal{G}\,\dCeff\bigr)^{-1}, \ }
\end{equation}
with $\dCeff$ the projected coupling of \S\ref{sec:testspace} and $\mathcal{G}$
the self propagator moment of the cube.

\subsection{What is new and what is borrowed}

New is $\dCeff$.  Borrowed is $\mathcal{G}$: the first-order route changes the
\emph{coupling}, not the Green's function.  The two moments it needs are
\begin{equation}\label{eq:selfmoments}
  \bG_{ij} = \int_V G^0_{ij}\,dV,
  \qquad
  M_{in,pk} = \int_V \partial'_p\partial'_k G^0_{in}\,dV .
\end{equation}

These are \emph{derived} in \S\ref{sec:moments} below, not taken from a table.
In the parameter basis
\begin{equation}\label{eq:Gmoment}
  \mathcal{G} = \frac{1}{V}
  \begin{pmatrix}\Gamma_0\,\bI_3 & \mathbf{0}\\ \mathbf{0} & -\bS\,\bW^{-1}\end{pmatrix},
  \qquad \bW = \operatorname{diag}(1,1,1,2,2,2),
\end{equation}
where $\bW$ undoes the engineering doubling of the shear slots.

Note that \eqref{eq:Gmoment} is asserted, not derived --- which is exactly why
\S\ref{sec:calibration} exists.

\subsection{The self moments, derived}\label{sec:moments}

The integrand of $M$ goes as $1/r^3$, which is not absolutely integrable in
three dimensions.  The customary remedy is to quote a tabulated Eshelby tensor
and add a delta-function term at $r=0$ by hand; omit the delta and the
depolarisation comes out with the \emph{wrong sign}.  Nothing below is
tabulated and nothing is added by hand.

\subsubsection{The static Green's tensor}

In the static limit the elastodynamic tensor reduces to Kelvin's,
\begin{equation}\label{eq:kelvin}
  G^0_{in}(\mathbf{r}) = \frac{a_0\,\delta_{in}}{r} + \frac{b_0\,x_ix_n}{r^3},
  \qquad
  a_0 = \frac{1}{8\pi}\Bigl(\frac1\mu + \frac{1}{\lambda+2\mu}\Bigr),
  \quad
  b_0 = \frac{1}{8\pi}\Bigl(\frac1\mu - \frac{1}{\lambda+2\mu}\Bigr),
\end{equation}
the coefficients being fixed by requiring that \eqref{eq:kelvin} satisfy the
Navier operator away from the origin.  An equivalent form makes the
distributional content visible, since $\partial_i\partial_n r = \delta_{in}/r -
x_ix_n/r^3$:
\begin{equation}\label{eq:kelvinhc}
  G^0_{in} = \frac{1}{4\pi\mu}\Bigl[\frac{\delta_{in}}{r}
             + h_c\,\partial_i\partial_n r\Bigr],
  \qquad h_c = -\frac{\lambda+\mu}{2(\lambda+2\mu)} ,
\end{equation}
and $a_0 = (1+h_c)/4\pi\mu$, $b_0 = -h_c/4\pi\mu$ identically.

\subsubsection{The moment as a distribution, and why no delta is inserted}

Write the moment as a pairing against the test function $w\,\mathbf{1}_V$
rather than as an integral of a function.  With $T = \partial_{\text{rest}}r^m$
and $\partial_q\mathbf{1}_V = -n_q\,\delta_{\partial V}$,
\begin{equation}\label{eq:recursion}
  \langle \partial_q T,\; w\,\mathbf{1}_V\rangle
  = \oint_{\partial V} n_q\,w\,T\,dA
  \;-\; \langle T,\; (\partial_q w)\,\mathbf{1}_V\rangle .
\end{equation}
That is the entire construction.  Two features make it exact rather than merely
plausible:

\begin{itemize}
\item $T$ is only ever evaluated \emph{on the boundary}, at distance $\Delta/2$
      from the origin, where every derivative of $r^m$ is smooth and bounded.
      The singular pointwise structure of $\partial^Dr^m$ at the origin is never
      formed, so there is no principal value to take and no delta to insert.
\item the recursion strictly lowers both the derivative count and the weight
      degree, terminating at $\int_V x_w\,r^m\,dV$, an ordinary convergent
      integral for $m\ge-1$.
\end{itemize}

The delta term is therefore neither added nor droppable --- it is already inside
the surface integral.  The sharp test is the Laplacian sum rule
$\sum_p E[-1;\{p,p\}] = -4\pi$, which a spherical excision returns as $0$.

\subsubsection{The face integrals, and one of them by pure geometry}

On the face $x_3 = \Delta/2$, scaling the in-face coordinates by $\Delta/2$
reduces every surface term to one of three integrals over $[-1,1]^2$:
\begin{equation}\label{eq:faceints}
  j_1 = \!\iint\!\frac{du\,dv}{(1+u^2+v^2)^{3/2}} = \frac{2\pi}{3},
  \quad
  j_2 = \!\iint\!\frac{u^2\,du\,dv}{(1+u^2+v^2)^{5/2}} = -\frac{2(\sqrt3-\pi)}{9},
  \quad
  k_1 = \!\iint\!\frac{du\,dv}{(1+u^2+v^2)^{5/2}} = \frac{2(2\sqrt3+\pi)}{9}.
\end{equation}
All three are evaluated in closed form.  The first needs no integration at all:
$j_1$ \emph{is} the solid angle subtended by one face at the cube's centre, and
the six faces subtend $4\pi$, so $j_1 = 4\pi/6 = 2\pi/3$.  That identity is a
check on the scaling and orientation of \eqref{eq:faceints}, obtained
independently of any quadrature.

The zeroth moment needs one further integral, which also has a closed form:
\begin{equation}\label{eq:ivr}
  \int_{[-1,1]^3}\frac{dV}{r}
  = -2\bigl(\pi + \log 64 - 12\log(1+\sqrt3)\bigr)
  = \tfrac43\log\bigl(70226 + 40545\sqrt3\bigr) - 2\pi
  = 9.520309455918214\ldots
\end{equation}
The two closed forms are independent evaluations of the same constant and agree
to the precision carried; the second is the one this package uses elsewhere.

\subsubsection{The closed forms}

Assembling \eqref{eq:kelvin} through the recursion \eqref{eq:recursion} with the
face integrals \eqref{eq:faceints} gives, with
\begin{equation}\label{eq:Mdecomp}
  M_{in,pk} = A\,\delta_{in}\delta_{pk}
            + B\,(\delta_{ip}\delta_{nk} + \delta_{ik}\delta_{np})
            + C\,E_{inpk},
  \qquad
  E_{inpk} = \begin{cases}1 & i=n=p=k\\ 0 & \text{otherwise,}\end{cases}
\end{equation}
the third structure being the one a cubic --- as opposed to isotropic ---
geometry admits,
\begin{align}
  B &= \frac{\sqrt3\,(\lambda+\mu)}{6\pi\mu(\lambda+2\mu)}, \label{eq:Bclosed}\\
  A &= B - \frac{1}{3\mu}, \label{eq:Aclosed}\\
  C &= -\Bigl(5 - \frac{2\pi}{\sqrt3}\Bigr)B, \label{eq:Cclosed}\\
  G_{ij} &= \delta_{ij}\,\Bigl(a_0 + \frac{b_0}{3}\Bigr)\Bigl(\frac{\Delta}{2}\Bigr)^{\!2}
            \int_{[-1,1]^3}\frac{dV}{r}. \label{eq:Gclosed}
\end{align}
$A$, $B$ and $C$ are scale-free; $G$ carries $\Delta^2$.  Note that $A$, $B$ and
$C$ are \emph{not} three independent constants --- the whole rank-four moment is
carried by $B$ together with the elementary $1/3\mu$, which is a structural
statement the tabulated presentation obscures.  $G_{ij}$ must be a multiple of
$\delta_{ij}$, because a rank-two cubic invariant has no other form: the cube
cannot distinguish its own axes at that order.

\subsubsection{The check that certifies the delta content}

From \eqref{eq:kelvinhc} with $\nabla^2(1/r) = -4\pi\delta^3$ and
$\nabla^2(\partial_i\partial_n r) = 2\partial_i\partial_n(1/r)$, whose cube
integral is $-\tfrac{4\pi}{3}\delta_{in}$,
\begin{equation}\label{eq:trace}
  \sum_p M_{in,pp}
  = -\,\frac{\delta_{in}}{\mu}\Bigl(1 + \frac{2h_c}{3}\Bigr)
  = -\,\delta_{in}\,\frac{2\lambda+5\mu}{3\mu(\lambda+2\mu)} ,
\end{equation}
derived without reference to the recursion at all.  The assembled moment
satisfies \eqref{eq:trace} identically, and the export refuses to write its
output otherwise.  A spherically excised moment returns zero on the left-hand
side, so \eqref{eq:trace} is exactly the statement that the $r=0$ content is
present and correctly weighted.

\subsection{The tabulated route agrees}

This package also carries the conventional tabulated Eshelby route, in which
$j_1$, $j_2$ and $k_1$ are quoted rather than integrated.  Evaluated alongside
\eqref{eq:Bclosed}--\eqref{eq:Gclosed} at $\lambda = 17.5$, $\mu = 22.5$ GPa,
$\Delta = 1$:

\begin{center}
\begin{tabular}{lrrr}
\toprule
moment & integrated & tabulated & relative\\
\midrule
$A$ & $-1.220110745874\times10^{-11}$ & $-1.220110745874\times10^{-11}$ & $0$\\
$B$ & $\phantom{-}2.613707356074\times10^{-12}$ & $\phantom{-}2.613707356074\times10^{-12}$ & $3.1\times10^{-16}$\\
$C$ & $-3.587055298887\times10^{-12}$ & $-3.587055298887\times10^{-12}$ & $2.3\times10^{-16}$\\
$G$ & $\phantom{-}6.622004020056\times10^{-12}$ & $\phantom{-}6.622004020056\times10^{-12}$ & $0$\\
\bottomrule
\end{tabular}
\captionof{table}{The table reproduces the integration to machine precision.
The direction of that statement matters: the integration \emph{verifies} the
table; the assembly reads the integrated values.}
\end{center}

The dynamic (smooth radiation) part carries no table --- it is already exact
polynomial integration --- and is taken from the existing route with its
tabulated static piece removed and the integrated one substituted.

\subsection{Two implementations of the coupling}

$\dCeff$ is rebuilt in Python directly from $\dA$, independently of the symbolic
derivation, and compared against the values that derivation exports.  Worst
relative difference over three contrasts: $1.4\times10^{-15}$ for $\dCeff$ and
exactly zero for the Navier coupling.  Two implementations agreeing is the
evidence standard; one agreeing with itself is not.

\subsection{Calibration: every normalisation fixed against existing code}
\label{sec:calibration}

Run the \emph{same} assembly with the ordinary Navier coupling in place of
$\dCeff$.  It must reproduce this package's validated four amplification
factors.  That pins the volume factors, the sign of $\mathcal{G}$ and the
weight $\bW$ --- the \emph{bookkeeping} of \eqref{eq:Gmoment} --- with no
freedom left to tune.  It does not certify the moments themselves; that is what
the integration of the previous subsection is for.  Against
$\alpha = 5000$, $\beta = 3000$, $\rho = 2500$ (SI), $\omega = 60$, $a = 0.5$:

\begin{center}
\begin{tabular}{llrr}
\toprule
case & factor & assembled vs existing (rel.) \\
\midrule
pure shear $(\Delta\mu = 1$ GPa$)$ & $A_u$, $A_\theta$, $A_e^{\text{diag}}$, $A_e^{\text{off}}$
  & $0$, $2.2\times10^{-16}$, $0$, $0$\\
moderate $(2,\,1,\,100)$ & same four
  & $0$, $2.3\times10^{-16}$, $1.1\times10^{-16}$, $0$\\
\bottomrule
\end{tabular}
\captionof{table}{The assembly reproduces the existing amplification factors to
machine precision.  Only after this is the coupling swapped.}
\end{center}

\paragraph{One term had to be isolated to get there.}  The weak form's
$\int_V\Delta\rho\,\omega^2\,\boldsymbol\phi\cdot\bu\,dV$ contributes to the
\emph{strain}--strain block as well as the rigid one, because an affine trial
displacement is not constant: $\int_V x_ax_b\,dV = \delta_{ab}Va^2/3$.  This is
a finite-size dynamic coupling between the density contrast and the strain
channels, and the factorised amplification factors do not carry it --- there
$A_u$ and the strain amplifications are independent.  It shifts $A_\theta$ by
$1.5\times10^{-7}$ at the moderate contrast and vanishes identically with
$\Delta\rho$.  It is a small correction that the factorisation drops, not a
defect in either object, and it is reported rather than absorbed.

\subsection{The first-order \texorpdfstring{$T$}{T}, channel by channel}

A single norm would average a correction that lives in a few channels against
others that do not move.  The comparison is therefore resolved by channel.

\begin{center}
\begin{tabular}{llrr}
\toprule
case & channel & relative difference from the Navier route\\
\midrule
\multirow{5}{*}{pure shear}
 & $\theta$ (bulk)            & $1.13\times10^{-4}$\\
 & $e_{\text{diag}}$          & $1.1\times10^{-16}$\\
 & $e_{12}$ (in-plane shear)  & $0$\\
 & $e_{13}$ (out-of-plane shear) & $\mathbf{8.36\times10^{-4}}$\\
 & $u$ (density)              & $0$\\
\midrule
\multirow{5}{*}{moderate}
 & $\theta$ (bulk)            & $1.74\times10^{-3}$\\
 & $e_{\text{diag}}$          & $1.1\times10^{-16}$\\
 & $e_{12}$                   & $0$\\
 & $e_{13}$                   & $\mathbf{8.36\times10^{-4}}$\\
 & $u$                        & $0$\\
\bottomrule
\end{tabular}
\captionof{table}{Amplification $(\bI - \mathcal{G}\Delta\mathbf{C})^{-1}$ along
each channel.  The density channel is identical in both routes, as
\S\ref{sec:whatitbuys} requires; the in-plane shear is untouched; the bulk and
out-of-plane shear channels move.}
\end{center}

\subsection{The broken symmetry, now visible in \texorpdfstring{$T$}{T} itself}

$e_{12}$ and $e_{13}$ are equivalent under $\Oh$.  For an isotropic cube in an
isotropic background their amplifications \emph{must} be equal.  They are, in the
Navier route, to exactly zero.  In the first-order route they split:

\begin{center}
\begin{tabular}{lrr}
\toprule
 & Navier & first-order\\
\midrule
$|A(e_{12}) - A(e_{13})|\,/\,|A(e_{12})|$ & $0$ & $8.36\times10^{-4}$\\
\bottomrule
\end{tabular}
\captionof{table}{The $\Oh$ violation of \S\ref{sec:oh}, measured on the
$T$-matrix rather than on the coupling.  Identical at both contrasts, because
both carry the same $\Delta\mu$ and the shear channels see only that.}
\end{center}

This is the sharpest form of the result.  It is not a norm over a matrix: it is
two numbers that a symmetry argument says must agree, and do not.

\paragraph{What the split does and does not license.}  It is tempting to read
this as the first-order $T$ failing an external test while the ordinary route
passes.  That reading claims more than the check can deliver, and the two halves
must be kept apart.

\emph{What it gives.}  An \textbf{arbiter-free lower bound on the error of the
first-order $T$}.  If $A(e_{12})$ and $A(e_{13})$ must be equal and differ by
$\delta$, then whichever lies closer to the truth, the larger of the two errors
is at least $\delta/2 \approx 4.2\times10^{-4}$.  An absolute error estimate
obtained from a symmetry alone, with nothing external to compare against, is
worth having --- and it is the absolute companion this project requires
alongside any ratio test.

\emph{What it does not give.}  Any bound whatever on the ordinary route's error.
A symmetric approximation can be symmetrically wrong, missing the truth in
\emph{both} channels by more than $8.4\times10^{-4}$ while satisfying the
symmetry exactly.  Symmetry compliance is not accuracy.  Setting a bound on one
formulation against no bound on the other is not a comparison, and no ordering
of the two follows from this table.

\emph{What is nonetheless attributable.}  The breaking is not manufactured by
the projection mismatch of \S\ref{sec:tmatrix}: $\mathcal{G}$ is itself
$\Oh$-symmetric, and \S\ref{sec:oh} measured the coupling's asymmetry
symbolically with no propagator present at all.  The asymmetry is a property of
the coupling, and it survives into $\mathbf{T}$.

\paragraph{The channel set has to include $e_{13}$.}  The coupling leaves
$e_{12}$ untouched by construction, so a comparison spanning only $\theta$,
$e_{\text{diag}}$ and $e_{12}$ reports that the shear channels agree --- true,
and uninformative.  A channel-resolved comparison is only as good as the
channels chosen, and the $\Oh$-equivalent pair is the one that carries the
information.

\subsection{What this does and does not settle}

The Born limits agree: at a contrast of $10^{-6}$ the two couplings differ by
$4.2\times10^{-8}$, i.e.\ at $O(\Delta c)$ as they must.

\begin{quote}
\textbf{The caveat that no number above reveals.}  The propagator moment used is
the \emph{single} (collocation) one, while the $\Jsix$ pairing makes the scheme
Galerkin.  The consistent pairing needs the \emph{double} moment.  This is the
same class of mismatch the package has already measured as a genuine error
elsewhere, so the table above should be read as the first-order coupling
evaluated through a collocation propagator, not as the fully consistent
first-order $T$.  No winner is declared, and none should be read in.
\end{quote}

%=====================================================================
\section{The propagator moment from the first-order system}\label{sec:propmom}

\S\ref{sec:tmatrix} took $\mathcal{G}$ from the ordinary Kelvin tensor, so the
$T$-matrix assembled there is first-order in its \emph{coupling} only.  This
section replaces that half.  The plan is
\texttt{docs/plans/2026-09-18-first-order-propagator-moment.md}; what follows is
written as each step lands, and says plainly where the front is.

The object is the sandwiched form
\begin{equation}\label{eq:sandwich}
  \bigl[\dCeff\,\mathcal{G}\,\dCeff\bigr]_{kl}
  = \bigl\langle \Jsix\psi_k \bigm| \dA\,\boldsymbol\Gamma\,\dA \bigm| \psi_l \bigr\rangle .
\end{equation}
Plain Galerkin on $\bq$ is not available: its mass matrix
$\langle\Jsix\psi_k,\psi_l\rangle$ is antisymmetric, and an antisymmetric matrix
in \emph{odd} dimension --- here nine --- is singular.

\subsection{Step 1: \texorpdfstring{$\Gamma$}{Gamma}, built from
\texorpdfstring{$\Aop$}{A} alone}

For $\partial_3\bq = \bA\bq + \bd\,\delta(z-z')$ the radiating Green's matrix is
\begin{equation}\label{eq:gammaspec}
  \boldsymbol\Gamma(z,z') =
  \begin{cases}
    +\bP_d\,e^{\bA(z-z')}, & z > z',\\[2pt]
    -\bP_u\,e^{\bA(z-z')}, & z < z',
  \end{cases}
\end{equation}
with $\bP_d,\bP_u$ the spectral projectors onto the down- and up-going
eigenspaces, so that $\boldsymbol\Gamma$ decays away from the source in both
directions and jumps by $\bP_d + \bP_u = \bI$ across it.

\paragraph{Projectors, not eigenvectors.}  The thesis writes the eigen-matrix
column by column with normalisation factors $\epsilon_P,\epsilon_S,\epsilon_H$
left symbolic.  Those factors are fixed by a flux normalisation that degenerates
at normal incidence, where the horizontal direction no longer distinguishes the
$SV$ and $SH$ polarisations --- exactly the $k_\parallel\to0$ difficulty.  The
spectral projector
\begin{equation}\label{eq:projector}
  \bP_d = \bV_d\bigl(\bW_d^{t}\bV_d\bigr)^{-1}\bW_d^{t}
\end{equation}
is invariant under any rescaling of the columns of $\bV_d$, so building
$\boldsymbol\Gamma$ from projectors removes the ambiguity instead of carrying
it.  That invariance is verified directly, by rescaling every eigenvector by a
random complex factor and confirming $\bP_d$ is unchanged.

\paragraph{The matrix exponential must never be formed.}
Equation~\eqref{eq:gammaspec} is a statement, not an algorithm.  Evaluating it as
written --- projecting \emph{after} exponentiating --- builds the growing branch
and then cancels it.  At a separation of $640$\,km the growing entries reach
$e^{+28}\approx10^{12}$, so the cancellation leaves $10^{-4}$ of absolute noise
on a true answer of $3\times10^{-8}$.  This is the Thomson--Haskell instability,
the reason invariant imbedding exists, reappearing inside our own construction;
it announces itself as a decay rate that turns \emph{negative} at large
separation.  Applying the exponential to the projected eigenvalues never forms
the growing branch at all,
\begin{equation}\label{eq:gammastable}
  \boldsymbol\Gamma(z-z') = \pm\!\!\sum_{i\,\in\,\text{down/up}}\!\!
    e^{s_i(z-z')}\,\mathbf{v}_i\mathbf{w}_i^{t},
  \qquad \mathbf{w}_i\!\cdot\!\mathbf{v}_j = \delta_{ij},
\end{equation}
and is what is used.  The moment integral needs only $|z-z'|\le 2a$, where the
instability cannot bite; a layered background with thick layers is where it
would.

\paragraph{$\omega$ is complex from the outset.}  The vertical wavenumbers
$k_{z,c} = \sqrt{k_c^2-k^2}$ vanish on the circles $|k| = k_P$ and $|k| = k_S$,
which the lateral integral must cross, and $k_z$ sits in a denominator.  At real
$\omega$ a quadrature straddles an integrable $1/\sqrt{\cdot}$ singularity and
does not converge; the directional-sweep work measured $1.9\times10^{-1} \to
1.1\times10^{-1}$ over an eightfold refinement, non-monotone.  With damping
$0.02$ the smallest $|s|$ \emph{on} the $P$ circle is $0.240$ rather than zero.

\paragraph{Checks.}  Thirteen, all passing: the spectrum is
$\{\pm ik_{z,P}\}$ and $\{\pm ik_{z,S}\}$ twice; $\bP_d+\bP_u=\bI$, both
idempotent, both commuting with $\bA$; $\bP_d$ invariant under rescaling;
$\partial_z\boldsymbol\Gamma = \bA\boldsymbol\Gamma$ on both sides; the jump
exactly $\bI$; and the decay rate converging to the slowest eigenvalue,
$0.034990 \to 0.033901 \to 0.032857 \to 0.032436$ against $0.032393$ --- a rate,
not a bar, because at any one finite separation the faster mode still
contributes.

\subsection{Step 2: the vertical integral, in closed form}

The vertical integral is exact, but it is \emph{not} an integral over the
square.  The two halves of $\Gamma$ carry different matrices, so the double
integral splits over the two triangles and they cannot be recombined: treating
the kernel as a function of $|z-z'|$ alone would symmetrise the propagator and
destroy the up/down distinction the first-order system exists to carry.

With affine trial fields the weights are $1$ and $z$, so four triangular
integrals are needed,
\begin{equation}\label{eq:tridef}
  T_{mn}(q) = \int_{-a}^{a}\!\!dz\int_{-a}^{z}\!\!dz'\;
              z^{m}z'^{n}\,e^{-q(z-z')},
\end{equation}
the lower triangle being $T_{nm}$ by relabelling.  All four are elementary:
\begin{align}
  T_{00} &= \frac{2aq - 1 + e^{-2aq}}{q^{2}}, \label{eq:T00}\\
  T_{10} &= -T_{01} = \frac{aq - 1 + (1+aq)\,e^{-2aq}}{q^{3}}, \label{eq:T10}\\
  T_{11} &= \frac{3 - 3(1+aq)^{2}e^{-2aq} + a^{2}q^{2}(2aq-3)}{3q^{4}}.
           \label{eq:T11}
\end{align}
Summing the two triangles with a common matrix recovers the familiar square
integral $2(2aq-1+e^{-2aq})/q^{2}$, and the odd weight cancels between them by
parity --- both checked, as is the large-$q$ behaviour $T_{00}\sim 2a/q$, each
triangle carrying half of the square's $4a/q$.

The decisive check of this step is not the closed forms themselves but the
\emph{assembly}: the vertical integration of $\Gamma$ performed through
\eqref{eq:T00}--\eqref{eq:T11}, matrices and all, reproduces a brute-force
double quadrature of $\Gamma$ to $2.5\times10^{-9}$ relative.

\subsection{Step 3: the lateral integral}

Count powers as $|k|\to\infty$ for the \emph{untransferred} form, in which
$\partial_\alpha\to ik_\alpha$: the box form factor gives $k^{-2}$,
\eqref{eq:T00} gives $q^{-1}\sim|k|^{-1}$, the two contrast operators give
$k^{2}$, and the measure gives $k$.  The strain channel does not converge.
(The transferred form of \S\ref{sec:transferred} does; that is one of the things
the augmentation buys, and it is taken up there.)

The divergence is not an artefact to be trimmed.  A delta function in real space
is a constant in $k$-space, so the non-decaying tail of the integrand \emph{is}
the Eshelby term of \S\ref{sec:moments}.  A bare cutoff discards precisely the
piece whose omission reverses the sign of the depolarisation.  Hence
split-and-subtract,
\begin{equation}\label{eq:splitsub}
  I(\mathbf{k}) = \bigl[I(\mathbf{k}) - I_\infty(\mathbf{k})\bigr] + I_\infty(\mathbf{k}),
\end{equation}
with the bracket integrated numerically --- convergent, and its convergence
measured as a rate --- and $\int I_\infty\,d^2k$ evaluated in closed form over
the whole plane, where the delta content resides.  The asymptote is taken of the
$z$-integrated kernel \eqref{eq:T00}--\eqref{eq:T11}, with the box form factors
kept exact: they are what makes the remainder converge, so expanding them would
defeat the subtraction.

\subsubsection{The scaling that makes the expansion uniform}

$\bA$ is not uniformly $O(k)$.  Rows $1$--$3$ carry $k$, but rows $5$--$6$ carry
$k^2$, because the basis mixes displacement with traction and the two have
different dimensions; expanding $\bA$ in $1/k$ as it stands mixes orders between
blocks.  The traction of an evanescent field scales as $\mu k u$, so the natural
balance is $\bq \to (\bu, \bta_3/k)$.  Under
$\bS = \operatorname{diag}(1,1,1,k,k,k)$,
\begin{equation}\label{eq:scaling}
  \widetilde{\bA} = \bS^{-1}\bA\,\bS
\end{equation}
is uniformly $O(k)$: the upper-right block gains a factor $k$, the lower-left
loses one, and the diagonal blocks are untouched.  Measured, $\|\bA\|/k$ grows
as $5.3\times10^{3} \to 5.3\times10^{5}$ over two decades in $k$ while
$\|\widetilde{\bA}\|/k$ settles at $53.0100 \to 53.0101$.  Being a similarity
transform it changes no eigenvalue, and $|s|/k \to 1$ for all six.

The scaling is also what makes the decomposition numerically usable.  On the
unscaled matrix the down- and up-going eigenvectors become nearly parallel as
$k$ grows --- the evanescent limit degenerates --- the eigen-decomposition is
reported badly conditioned, and the kernel entries collapse into rounding error,
from which no asymptotic power can be read.  The decomposition is therefore done
on $\widetilde{\bA}$ and mapped back, $\bK = \bS\,\widetilde{\bK}\,\bS^{-1}$,
checked against the verified kernel of step 2 at moderate $k$.

\subsubsection{Which entries fail to decay}

The order in $1/k$ that each entry of the $z$-integrated kernel carries is read
off numerically, over decades in $k$, rather than derived: a measured power is
harder to fool oneself with than one obtained from a $6\times6$
eigen-perturbation.  A dot marks an entry that vanishes, so no power can be read.

\begin{center}
\begin{tabular}{l|cccccc}
\toprule
 & $u_z$ & $u_x$ & $u_y$ & $\tau_{zz}$ & $\tau_{xz}$ & $\tau_{yz}$\\
\midrule
$u_z$        & $\cdot$ & $-1$ & $-1$ & $-2$ & $\cdot$ & $\cdot$\\
$u_x$        & $-1$ & $\cdot$ & $\cdot$ & $\cdot$ & $-2$ & $-2$\\
$u_y$        & $-1$ & $\cdot$ & $\cdot$ & $\cdot$ & $-2$ & $-2$\\
$\tau_{zz}$  & $\mathbf{0}$ & $\cdot$ & $\cdot$ & $\cdot$ & $-1$ & $-1$\\
$\tau_{xz}$  & $\cdot$ & $\mathbf{0}$ & $\mathbf{0}$ & $-1$ & $\cdot$ & $\cdot$\\
$\tau_{yz}$  & $\cdot$ & $\mathbf{0}$ & $\mathbf{0}$ & $-1$ & $\cdot$ & $\cdot$\\
\bottomrule
\end{tabular}
\captionof{table}{Power of $k$ carried by each entry of the $z$-integrated
kernel as $|k|\to\infty$.  Exactly one block fails to decay.}
\end{center}

\begin{quote}
\textbf{The entire ultraviolet problem sits in one block:} the traction rows
against the displacement columns.  Everything else decays at least as $1/k$ and
needs no subtraction.
\end{quote}

That block is also the one the physics predicts.  In the evanescent limit a
short-wavelength displacement produces a traction $\tau\sim\mu k u$, while the
vertical integral \eqref{eq:T00} supplies $1/k$; the ratio is $k$-independent.
A traction response to a displacement that does not depend on $k$ \emph{is} a
local stiffness --- which is the Eshelby delta of \S\ref{sec:moments}, seen from
the first-order side.  The two representations identify the same object, and
$I_\infty$ has to reproduce exactly this block.

\subsubsection{The asymptote in closed form}

The limit of the non-decaying block is not merely numerical.  Writing
$\hat{\mathbf{k}}$ for the lateral direction,
\begin{align}
  C\bigl[\tau_{zz}\!\leftarrow\!u_z\bigr] &= -\nu_1
     = -\frac{4\mu(\lambda+\mu)}{\lambda+2\mu}, \label{eq:Cnu1}\\
  C\bigl[\tau_{\alpha z}\!\leftarrow\!u_\beta\bigr]
     &= \mu\bigl(\hat{k}_\alpha\hat{k}_\beta - \delta_{\alpha\beta}\bigr).
     \label{eq:Cinplane}
\end{align}

\begin{quote}
\textbf{$\nu_1$ is not a new constant.}  It is the same $\nu_1$ that appears in
$\Aop_{12}$ --- born in the elimination \eqref{eq:U}, and shown in
\S\ref{sec:arbiterB} to be the thesis's $\zeta$.  The subtraction is expressible
entirely in the operator's own constants.  A wrongly identified asymptote would
have no reason to land on $\nu_1$ and $\mu$.
\end{quote}

Both are established as \emph{rates}.  The relative residuals fall
$1.50\times10^{-3} \to 1.50\times10^{-4} \to 1.50\times10^{-5}$ across decades in
$k$, ratios $10.0006$ and $10.011$ --- clean $O(1/k)$, which is the $O(k_c/k)$
correction to the evanescent limit.  Agreement at a single $k$ would not
distinguish a correct closed form from a coincidence.

The in-plane form is the sharper test.  Two constants are fitted on \emph{two}
angles and then predict \emph{six}: $P = -22.4977$ and $Q = +22.4948$, both
equal to $\mu = 22.5$, with predictions tracking measurements to five digits
across $\theta\in[0,2.4]$.  The isotropic entry \eqref{eq:Cnu1} is constant to
seven digits over the same range.

\paragraph{What this settles about the integration.}  The asymptote \emph{is}
direction-dependent --- the relative variation over $\theta$ is $0.39$ --- so
$I_\infty$ is not a pure delta and the regulator cannot be avoided.  But the
structure is the most favourable available: $\delta_{\alpha\beta}$ and
$\hat{k}_\alpha\hat{k}_\beta$ are the only two rank-two invariants a direction
can build, and they transform to a local delta and a principal-value dipole
kernel respectively, both standard.  The $\tau_{zz}\!\leftarrow\!u_z$ entry is
isotropic and so is a pure delta of weight $-\nu_1$.

\subsubsection{The lateral integral, and where it must be evaluated}

The box form factor is separable --- the lateral transform of the cube's
indicator is $(2a)^2\operatorname{sinc}(k_1a)\operatorname{sinc}(k_2a)$ --- so
with \eqref{eq:Cnu1}--\eqref{eq:Cinplane} and one lateral derivative from each
$\dA$, the integrand carries $k_ck_d$ against $C$, giving two families:
$\int d^2k\,W k_ck_d$, which is separable, and $\int d^2k\,W k_ak_bk_ck_d/k^2$,
which is not.  The elementary one-dimensional integrals are
\begin{equation}\label{eq:sincints}
  \int_{-\infty}^{\infty}\!\operatorname{sinc}^2(ka)\,dk = \frac{\pi}{a},
  \qquad
  \int_{-\infty}^{\infty}\!k^2\operatorname{sinc}^2(ka)\,e^{-\epsilon|k|}dk
  = \frac{4}{4a^2\epsilon + \epsilon^3}
  = \frac{1}{a^2\epsilon} - \frac{\epsilon}{4a^4} .
\end{equation}

\begin{quote}
\textbf{The second is a pure pole with no finite part.}  So
$\int I_\infty\,d^2k$ is not a number, and the split-and-subtract cannot be
closed in $k$-space alone: adding back a pure pole restores exactly the
divergence it was introduced to remove.
\end{quote}

The method is not wrong; the tail is being evaluated in the wrong
representation.  Forming $|\widehat{S}|^2$ first and then integrating term by
term destroys the structure that makes the constant part finite.  By Parseval,
\begin{equation}\label{eq:parseval}
  \int\!\!\int f(\mathbf{x})\,K(\mathbf{x}-\mathbf{x}')\,g(\mathbf{x}')
  = \int\!\frac{d^2k}{(2\pi)^2}\,\widehat{f}(-\mathbf{k})\,
      \widehat{K}(\mathbf{k})\,\widehat{g}(\mathbf{k}),
\end{equation}
so a $\widehat{K}$ tending to a constant $C$ contributes
$C\!\int_V\! f g\,d\mathbf{x}$ --- a \emph{local} term, finite, needing no
regulator at all.  This is verified directly in one dimension, both sides
returning $1/2a$.

The decomposition therefore splits by representation rather than by magnitude:

\begin{center}
\begin{tabular}{lll}
\toprule
piece & real-space object & evaluation\\
\midrule
isotropic, $-\nu_1$ and $P\delta_{\alpha\beta}$ & delta & local, elementary\\
traceless $\hat{k}_\alpha\hat{k}_\beta - \tfrac12\delta_{\alpha\beta}$
  & 2-D principal-value dipole & conditionally convergent\\
remainder $I - I_\infty$ & --- & numerically, in $k$\\
\bottomrule
\end{tabular}
\captionof{table}{Split-and-subtract survives, but $\int I_\infty$ is evaluated
in real space.  Only the remainder stays in $k$.}
\end{center}

The traceless part is of degree $-2$ with vanishing angular average --- exactly
the conditional convergence the real-space engine of \S\ref{sec:moments} handles
by peeling derivatives onto the cube faces.  The two representations meet again
here: the same distributional structure, in two dimensions rather than three.

\subsubsection{The principal-value piece vanishes}

In two dimensions $\nabla^2\log\rho = 2\pi\delta^2$, and
$\partial_\alpha\partial_\beta\log\rho = (\delta_{\alpha\beta} -
2\hat{n}_\alpha\hat{n}_\beta)/\rho^2$ away from the origin; taking the trace
fixes the delta weight, so that
\begin{equation}\label{eq:pvsplit}
  \partial_\alpha\partial_\beta\log\rho
  = \mathrm{PV}\!\left[\frac{\delta_{\alpha\beta}
      - 2\hat{n}_\alpha\hat{n}_\beta}{\rho^{2}}\right]
    + \pi\,\delta_{\alpha\beta}\,\delta^{2}(\mathbf{x}) .
\end{equation}
By translation invariance the double integral over the square collapses onto the
autocorrelation of its indicator, $w(\mathbf{u}) = (2a-|u_1|)(2a-|u_2|)$, exactly
--- removing two of the four dimensions --- and both derivatives move onto $w$,
which is legitimate because $w$ has compact support.

$w$ is piecewise bilinear, so $\partial_1\partial_1 w$ is a sum of line deltas.
It has three terms, not one: $w$ vanishes at $|u_1| = 2a$ but
$\partial_1 w$ does \emph{not} --- it jumps from $\mp(2a-|u_2|)$ to zero --- so
there are edge deltas at $u_1 = \pm2a$ alongside the interior one at $u_1 = 0$.
With them,
\begin{equation}\label{eq:Dab}
  D_{11} = \underbrace{3}_{\text{interior}} + \underbrace{(\pi-3)}_{\text{edges}}
         = \pi = \pi(2a)^2 \Big|_{a=1/2},
  \qquad D_{12} = 0 ,
\end{equation}
and the trace identity $\operatorname{tr}D = \langle\nabla^2\log\rho, w\rangle =
2\pi w(0) = 2\pi(2a)^2$ holds exactly.  Dropping the edge deltas gives $6$ in
place of $2\pi$, so that identity is what detects them.

\begin{quote}
\textbf{Hence $\langle \mathrm{PV}_{\alpha\beta}, w\rangle = D_{\alpha\beta} -
\pi\delta_{\alpha\beta}w(0) = 0$: the traceless dipole contributes nothing.}
\end{quote}

It could not have contributed.  $\langle\mathrm{PV}_{\alpha\beta},w\rangle$ is a
symmetric two-tensor invariant under the square's $D_4$ symmetry, and the only
such tensor is a multiple of $\delta_{\alpha\beta}$ --- the traceless part has
nowhere to live.  The conditional convergence is real, but for this cross-section
it integrates away.

\begin{quote}
\textbf{The whole of $\int I_\infty$ is therefore the local delta term.}  The
piece that looked hardest --- the conditionally convergent, shape-dependent
dipole --- is identically zero here, and what remains is elementary.
\end{quote}

\subsubsection{Why the transferred form converges}\label{sec:transferred}

In the untransferred form $\partial_\alpha \to ik_\alpha$, each contrast
operator contributes a factor of $k$, and the pole of \eqref{eq:sincints}
follows.  After the by-parts transfer of \eqref{eq:byparts} the derivatives act
on the \emph{polynomials, in real space}: $\partial_\alpha$ of an affine trial
function is a constant, whose lateral transform is again the box form factor.
No factor of $k$ survives, and the lateral integral converges.

\begin{quote}
\textbf{The augmentation therefore does two things at once:} it keeps the
indicator from being differentiated, and it makes the lateral integral
convergent.  Both are the same integration by parts.
\end{quote}

Convergence is nevertheless slow.  With $\widehat{K}\to C$ the truncation error
falls only as $1/\Lambda$, which is what the subtraction is for.

\subsubsection{The tail, in closed form}

The tail is elementary.  For the in-plane entry
$C = \mu(\hat{k}_1^2-1) = -\mu\,k_2^2/k^2$, so
\begin{equation}\label{eq:tail}
  \int\! d^2k\,W\,C = -\mu\!\int\! d^2k\,
    \operatorname{sinc}^2(k_1a)\operatorname{sinc}^2(k_2a)\,\frac{k_2^2}{k^2}
  = -\frac{\mu}{2}\Bigl(\frac{\pi}{a}\Bigr)^{\!2},
\end{equation}
because $k_1^2/k^2 + k_2^2/k^2 = 1$ and the weight is symmetric under
$k_1\!\leftrightarrow\!k_2$, so each piece is exactly half of
$\int\!d^2k\,W = (\pi/a)^2$.  No quadrature is required; direct integration
returns $19.7376$ against $(\pi/a)^2/2 = 19.7392$, the difference being the
quadrature's own error on a slowly converging integrand.

\subsubsection{The remainder, and the test that closes the construction}

Truncating the lateral integral at $|k| < \Lambda$ for the
$\tau_{xz}\!\leftarrow\!u_x$ entry, with and without the subtraction:

The split is only worth making if both halves give the same answer.  Truncating
at $|k| < \Lambda$ for the $\tau_{xz}\!\leftarrow\!u_x$ entry:

\begin{center}
\begin{tabular}{rrrr}
\toprule
$\Lambda$ & unsubtracted & remainder & remainder $+$ tail\\
\midrule
$80$  & $-318.2214$ & $118.83504$ & $-325.29716$\\
$160$ & $-321.6764$ & $118.82727$ & $-325.30493$\\
$320$ & $-323.4539$ & $118.82513$ & $-325.30707$\\
$640$ & $-324.3400$ & $118.82466$ & $-325.30754$\\
\bottomrule
\end{tabular}
\captionof{table}{The subtracted route is converged to six figures by
$\Lambda = 80$; the unsubtracted one is still $0.97$ short of the limit at
$\Lambda = 640$.}
\end{center}

The remainder's successive differences fall $7.8\times10^{-3} \to
2.1\times10^{-3} \to 4.7\times10^{-4}$, a ratio near $0.27$, against the
unsubtracted $0.50$ --- the latter being the $1/\Lambda$ decay the constant part
imposes.  Extrapolating both:
\begin{equation}\label{eq:closure}
  \underbrace{-325.226}_{\text{unsubtracted limit}}
  \;\;\text{versus}\;\;
  \underbrace{-325.308}_{\text{remainder} + \text{closed-form tail}},
  \qquad \text{relative difference } 2.5\times10^{-4},
\end{equation}
and that residue is dominated by the single-term extrapolation of the slowly
converging unsubtracted series, not by the construction.  At $\Lambda = 640$ the
unsubtracted error is $0.97$ and the subtracted one $5\times10^{-4}$: a factor
of about two thousand.

\begin{quote}
\textbf{Split-and-subtract is verified.}  The two routes to the same number
agree, the tail is closed in form by \eqref{eq:tail}, and the remainder converges
three orders of magnitude faster than the direct integral.
\end{quote}

\emph{Status of Task~1.}  The method is established end to end: $\Gamma$ from
spectral projectors; the triangular $z$ integrals and their assembly; the
scaling, the power table, and the asymptote $-\nu_1$,
$\mu(\hat{k}\hat{k}-\delta)$ verified as an $O(1/k)$ rate; $\int I_\infty$
decomposed by representation, its principal-value part vanishing by $D_4$
symmetry and its remaining part closed by \eqref{eq:tail}; and the closure test
\eqref{eq:closure}.  What remains is bookkeeping rather than method: assembling
all nine trial functions and the full contrast operator, and reading the shear
moment $B$ off the result for comparison with
$\sqrt3(\lambda+\mu)/6\pi\mu(\lambda+2\mu)$.

%=====================================================================
\section{Summary of checks}

\begin{center}
\begin{tabular}{llr}
\toprule
object & what is established & checks\\
\midrule
$\Aop$ & built from $c_{ijkl}$; matches the published isotropic blocks & 4\\
       & matches the thesis $\bA(k_x,k_y)$ under one similarity transform & 3\\
       & symmetry relations, spectrum, normal incidence & 6\\
       & stiffness matrices, $\bU_{\alpha l}$, $\nu_1$, $\nu_2$ & 9\\
       & source vector, negative control & 3\\
\midrule
$\dA$  & exact contrast operator; Navier body force on a generic field & 4\\
       & nonlinearity in the contrast, quantified & 3\\
       & augmentation: constant, reproduces the form, divergence exhibited & 4\\
       & $\Oh$/trial-space structure, density channel & 6\\
\midrule
$\dCeff$ & $\Jsix\dA$ symmetric; test space derived; controls & 5\\
         & $9\times9$ symmetric, rank $9$, rigid block, basis trap & 7\\
         & Born agreement, $O(\Delta c^2)$ difference, localisation & 9\\
         & $\Oh$ decomposition, positive and negative controls & 6\\
\midrule
$\boldsymbol\Gamma$ & propagates by $\exp(\bA\,dz)$ with all multiples present & 3\\
\midrule
moments & face integrals evaluated; $j_1$ also by solid angle & 4\\
        & $\int dV/r$ in two independent closed forms & 1\\
        & $B$, and $A = B - 1/3\mu$, $C = -(5-2\pi/\!\sqrt3)B$, $G$ & 4\\
        & Kelvin coefficients against the $h_c$ form & 1\\
        & the trace identity \eqref{eq:trace}, derived without the recursion & 1\\
        & the tabulated route reproduces all four & 1\\
\midrule
$\mathbf{T}$ & coupling rebuilt independently, matches the symbolic derivation & 2\\
             & calibration: reproduces the existing amplification factors & 1\\
             & the finite-size density--strain term, isolated and sized & 2\\
             & density channel identical; $\Oh$ split on the $e_{12}/e_{13}$ pair & 6\\
             & Born limits agree & 1\\
\bottomrule
\end{tabular}
\captionof{table}{All checks pass.  The symbolic scripts report $25/25$ and
$44/44$, the $T$-matrix assembly $12/12$, and the propagator gate exits clean.}
\end{center}

%=====================================================================
\section{What is not established}\label{sec:open}

In keeping with the project's standard that an unvalidated formula be labelled
as such, the following are \emph{open}:

\begin{enumerate}
\item \textbf{The $T$-matrix uses the collocation propagator, not the
      consistent one.}  $\mathbf{T}$ is assembled in \S\ref{sec:tmatrix} and
      calibrated exactly, but $\mathcal{G}$ there is the \emph{single} cube
      moment while the $\Jsix$ pairing makes the scheme Galerkin.  The
      consistent pairing needs the \emph{double} moment, which has not been
      extracted.  Everything in \S\ref{sec:tmatrix} is therefore the first-order
      coupling seen through a collocation propagator.
\item \textbf{No comparison against an independent arbiter has been run.}  The
      quantities in \S\ref{sec:whatitbuys} and \S\ref{sec:tmatrix} compare two
      \emph{formulations} with each other, not either against truth.  The $\Oh$
      split is the one partial exception, and it is one-sided: it lower-bounds
      the first-order error at $\approx4.2\times10^{-4}$ and bounds the ordinary
      route not at all.  \textbf{Which formulation is more accurate remains
      open}, and nothing in this note decides it.
\item \textbf{Whether the resummation improves accuracy is unknown.}  It is
      exact as a local inversion; that is not the same claim.
\item \textbf{$\Oh$ invariance was tested on a subgroup of order $3$}, not the
      full group.
\item \textbf{The derivative count for a collocation scheme differs from the
      Galerkin one.}  In the Galerkin pairing the transferred derivative lands
      on polynomials, so the third-derivative moment cannot arise; under
      collocation it transfers onto $\boldsymbol\Gamma$ and can.  Do not quote
      the Galerkin count for a collocation build.
\end{enumerate}

\section{Where the route belongs}

The $x_3$ preference of \S\ref{sec:oh} is an artefact only because the problem
posed there --- an isotropic cube in an isotropic whole space --- has no
preferred direction.  A layered background does.  There the formulation's
directional structure matches the physics rather than contradicting it, the
$\Oh$ argument does not apply, and the propagator is already available and
already bridged to $\Aop$ by \eqref{eq:bridge}.

That is also the case in which the ordinary route has no competitor: the cube
moment method rests on volume moments of the whole-space Green's tensor, and in
a stratified background that object does not exist.  The first-order form
supplies what the moment method cannot.

\end{document}
